% ====================================================================
%  Inderscience — International Journal of Biometrics (IJBM)
% --------------------------------------------------------------------
%  Template note: Inderscience's official submission system accepts
%  manuscripts in Microsoft Word (template:
%  inderscience.com/info/inauthors/author_journal.php) OR LaTeX. For
%  LaTeX, Inderscience asks authors to use their generic
%  `interscience_paper.cls` class on submission (or plain article and
%  the production team re-typesets). We use the plain `article` class
%  here so the draft compiles anywhere; swap the class at submission.
%  Inderscience expects: abstract (max 150 words), 4-6 keywords,
%  Harvard-style references, biographical notes for each author.
%  Target length: 6000-8000 words (12-18 typeset pages).
% ====================================================================
%% STATUS 2026-07-03: Experimental design superseded by the project audit
%% (research_logs/2026-07-03-project-audit-sota-roadmap.md). This venue awaits
%% the P4 flagship cross-lingual (Sinhala<->Tamil) study. Benchmark v0 results
%% now exist (see research_logs/2026-07-03-*) and are quoted below as
%% preliminary context; final tables require the SLCeleb v1 benchmark and the
%% P4 runs (cross-lingual trials, DANN/Dual-LoRA arms).
% ====================================================================
\documentclass[11pt,a4paper]{article}

% Engine guard: fontspec under XeLaTeX/LuaLaTeX (as in the 2026-05-23
% builds); T1/Latin Modern fallback so the draft also compiles with
% pdfLaTeX on machines without xetex/luaotfload.
\usepackage{iftex}
\iftutex
  \usepackage{fontspec}
  \setmainfont{DejaVu Serif}
  \setsansfont{DejaVu Sans}
  \setmonofont{DejaVu Sans Mono}
\else
  \usepackage[T1]{fontenc}
  \usepackage{lmodern}
\fi

\usepackage[a4paper,margin=1in]{geometry}
\usepackage[hidelinks]{hyperref}
\usepackage{booktabs,tabularx,amsmath,amssymb,microtype,enumitem,parskip}
\usepackage{authblk}

\title{Low-resource biometric speaker verification on Sinhala--Tamil
bilingual audio: a single-corpus ablation study with cross-lingual
transfer, score normalisation, and PLDA scoring}

\author[1]{[First Author]}
\author[1]{[Second Author]}
\affil[1]{[Department], [University], Sri Lanka. \\
\texttt{[email@institution]}}

\date{}

\begin{document}
\maketitle

\begin{abstract}
\noindent Voice-based biometric authentication has matured on English
benchmarks but remains under-evaluated for the low-resource languages
typical of South Asian deployment contexts. We present the first
reproducible speaker-verification (SV) benchmark for Sri Lankan Sinhala
and Tamil, evaluated on a two-tier protocol: benchmark v0, 527 speakers
of read speech from OpenSLR SLR52 and SLR65 (478 Sinhala, 49 Tamil;
12{,}000 trials per language), and benchmark v1, built on SLCeleb, our
280-speaker in-the-wild corpus (pending data placement).
Building on an ECAPA-TDNN encoder with AAM-Softmax training, we
systematically evaluate four authentication-relevant levers:
cross-lingual fine-tuning from an English-pretrained checkpoint with
selective front-end freezing; layer-wise learning-rate decay; an
auxiliary language-identification head trained jointly with the
speaker loss; and a two-stage PLDA followed by adaptive symmetric
score normalisation. Each lever is ablated independently, with three-
seed reproducibility. We report monolingual Sinhala-only, Tamil-only,
and cross-lingual equal-error rates, together with an open-source
toolkit and corpus-preparation pipeline released for downstream use.
\end{abstract}

\noindent \textbf{Keywords:} speaker verification; biometrics;
cross-lingual transfer; low-resource speech; Sinhala; Tamil; PLDA;
score normalisation; ECAPA-TDNN.

\section{Introduction}
\label{sec:intro}

Voice-based biometrics is increasingly deployed in financial
authentication, call-centre fraud detection, and forensic audio
analysis. The dominant published benchmarks
\cite{nagrani2017voxceleb} use English speech almost exclusively, and
the state-of-the-art architectures and training recipes
\cite{desplanques2020ecapa,deng2019arcface} were validated at that
scale. For Sri Lankan deployment contexts the linguistic mismatch is
fundamental: Sinhala (Indo-Aryan, [N=17] million speakers) and Tamil
(Dravidian, [N=5] million Sri Lankan speakers) are absent from the
publicly available SV corpora, bilingual and code-switched speech is
normal in everyday usage, and the labelled-speaker base accessible to
a single research effort is small (tens to low hundreds).

This paper contributes:
\begin{enumerate}
    \item A reproducible two-tier Sinhala/Tamil SV benchmark ---
    v0: 527 speakers (478 Sinhala from OpenSLR SLR52, 49 Tamil from
    OpenSLR SLR65 read speech); v1: SLCeleb, our 280-speaker
    in-the-wild corpus (pending) --- with monolingual, cross-lingual,
    and pooled trial lists.
    \item Single-feature ablations isolating the contribution of
    cross-lingual fine-tuning, layer-wise LR decay, multi-task
    language-ID, AS-Norm score normalisation, and PLDA scoring.
    \item An open-source toolkit
    (\texttt{voxceleb\_trainer-SL-SPV}) \cite{slspv2026repo}
    extending the upstream voxceleb\_trainer with eleven features and
    twenty-six bug fixes, each selectable via a single YAML key.
\end{enumerate}

The remainder of this paper is organised as follows.
Section~\ref{sec:related} reviews relevant work in SV architectures,
cross-lingual transfer, and score normalisation.
Section~\ref{sec:method} describes the corpus, architectures, and
training recipe.
Section~\ref{sec:experiments} reports results.
Section~\ref{sec:discussion} discusses biometric-deployment
implications, and Section~\ref{sec:conclusion} concludes.

\section{Related work}
\label{sec:related}

\subsection{Speaker verification architectures}
\label{sec:related:arch}

Modern SV pipelines extract a fixed-length embedding from variable-
length audio and score trial pairs by a similarity function on those
embeddings. ECAPA-TDNN \cite{desplanques2020ecapa} is the de-facto
reference architecture, combining Res2Net-style multi-scale 1D
convolutions, squeeze-and-excite channel re-weighting, multi-layer
feature aggregation, and channel-attentive statistical pooling. We
adopt ECAPA-TDNN-Large ($C=1024$ channels) as the headline encoder.

\subsection{Cross-lingual transfer for biometrics}
\label{sec:related:transfer}

When the target-language SV corpus is small, training from random
initialisation overfits. Cross-lingual fine-tuning from a
high-resource (typically English VoxCeleb1) checkpoint is the
established remedy. The fine-tune policy hyper-parameters — which
layers to freeze, the global LR multiplier, and whether to apply
layer-wise LR decay (LLRD) — are well studied in the BERT / WavLM
literature but less systematically reported for SV fine-tuning, which
motivates our explicit ablation in Section~\ref{sec:experiments}.

\subsection{Score normalisation and generative backends}
\label{sec:related:scoring}

Adaptive symmetric score normalisation (AS-Norm)
\cite{matejka2017analysis} computes per-trial-side cohort statistics
from a fixed set of impostor embeddings and normalises trial scores
against those statistics. Two-covariance PLDA
\cite{sizov2014unifying} replaces cosine scoring with a likelihood
ratio test under same-speaker vs different-speakers hypotheses. Both
are eval-time post-processing steps that do not require retraining,
and they compose: the standard NIST SRE pipeline order is
PLDA $\to$ AS-Norm.

\subsection{Auxiliary objectives for low-resource SV}
\label{sec:related:aux}

Multi-task language-ID heads — a linear classifier trained jointly
with the speaker loss to predict the utterance's language — regularise
the speaker embedding to retain language-discriminative structure
that aids cross-lingual discrimination. The opposite
(adversarial / DANN-style \cite{ganin2015dann}) is also studied but
empirically more brittle on small corpora; we therefore evaluate the
multi-task variant as the primary lever and treat DANN as
infrastructure-ready future work.

\section{Methodology}
\label{sec:method}

\subsection{Corpus}
\label{sec:method:corpus}

Evaluation is two-tier. Benchmark v0 comprises 527 speakers: 478
Sinhala from OpenSLR SLR52 (crowdsourced read speech, 16~kHz) and 49
Tamil from OpenSLR SLR65 (studio read speech, resampled to 16~kHz at
load), keeping speakers with at least 10 utterances. Benchmark v1 is
built on SLCeleb, our 280-speaker in-the-wild Sinhala/Tamil corpus
(multi-genre, multi-session, with bilingual speakers), pending data
placement. Within each speaker we hold out 20\% of utterances for
trial-pair construction; the remainder forms the training pool
(31{,}049 training utterances at a 60-utterance-per-speaker cap for
the fine-tuning experiments).

v0 trial pairs are constructed at a 1:3 target-to-impostor ratio at
three operating points: monolingual Sinhala (3{,}000 target +
9{,}000 impostor pairs), monolingual Tamil (3{,}000 + 9{,}000), and
cross-lingual (pairs for speakers present in both languages --- empty
at v0, since the OpenSLR corpora contain no bilingual speakers, and
populated at v1 by SLCeleb). Pair sampling is reproducible
from a single fixed random seed \cite{slspv2026repo}.

\subsection{Architecture and training}
\label{sec:method:arch}

The speaker encoder is ECAPA-TDNN \cite{desplanques2020ecapa} with
$C=1024$ channels, 80 log-mel filterbanks, embedding dimension 192,
and channel-attentive statistical pooling. The training loss is
additive-angular-margin softmax (AAM-Softmax) \cite{deng2019arcface}
with margin 0.2 and scale 30. Augmentation reuses the MUSAN +
RIR simulation chain at training time.

\subsection{Cross-lingual fine-tune protocol}
\label{sec:method:finetune}

For configurations P1 and P2 (Section~\ref{sec:experiments}) we
initialise from an ECAPA-TDNN checkpoint pretrained on VoxCeleb1
\cite{nagrani2017voxceleb}. During SL fine-tuning the mel front-end
is frozen, the global LR is scaled by $0.1\times$, and layer-wise
LR decay with $\gamma=0.9$ is applied across the four ECAPA layers
so that the layer closest to the loss head receives the full LR and
the input-side layer receives $0.9^3 \approx 0.73\times$ of it.

\subsection{Multi-task language-ID head}
\label{sec:method:langaux}

A linear head $W_{\mathrm{aux}}\in\mathbb{R}^{4\times 192}$ predicts
the utterance's primary language (si / ta / en / mix) from the speaker
embedding before the $n_{\mathrm{PerSpeaker}}$ reshape. The combined
loss is
\begin{equation}
\mathcal{L}_{\mathrm{total}}
= \mathcal{L}_{\mathrm{speaker}}
+ \lambda \, \mathcal{L}_{\mathrm{lang}},
\quad \lambda=0.3.
\end{equation}
Speakers absent from the per-speaker language lookup carry
\texttt{ignore\_index=-1} and contribute no gradient.

\subsection{Eval-time scoring backends}
\label{sec:method:scoring}

We evaluate three eval-time scoring backends:
\begin{enumerate}
    \item Cosine similarity on length-normalised embeddings (baseline).
    \item Two-covariance PLDA \cite{sizov2014unifying} with LDA
    pre-reduction to $d=200$ and length normalisation.
    \item AS-Norm \cite{matejka2017analysis} with top-$K=300$ cohort
    statistics on a 500-speaker in-domain cohort, layered on top
    of either backend (1) or (2).
\end{enumerate}

\subsection{Reproducibility}
\label{sec:method:repro}

All experiments use deterministic-mode training (cuDNN deterministic
flag, TF32 disabled,
\texttt{torch.use\_deterministic\_algorithms} enabled). Primary
configurations are run with three random seeds (42, 123, 7); we
report $\mathrm{mean} \pm \mathrm{std}$. Single-feature ablations
use seed 42 only.

\section{Experiments}
\label{sec:experiments}

\subsection{Configurations}
\label{sec:exp:configs}

We define three primary configurations and five single-feature
ablations:

\begin{itemize}[leftmargin=*]
    \item \textbf{P0:} ECAPA-TDNN from scratch on SL-SPV.
    \item \textbf{P1:} P0 + cross-lingual fine-tune (front-end frozen,
    LLRD).
    \item \textbf{P2:} P1 + multi-task language-ID head + AS-Norm +
    PLDA + per-language evaluation.
    \item Ablations of P2 with each of the five features individually
    disabled.
\end{itemize}

\subsection{Preliminary monolingual baselines (benchmark v0)}
\label{sec:exp:prelim}

%% Numbers below are from the 2026-07-03 benchmark v0 runs
%% (research_logs/2026-07-03-{zeroshot-baseline-table,p2-backend-adaptation,
%% p3-peft-finetuning}.md). Monolingual context only; the P4 cross-lingual
%% arms this paper targets are pending.

\begin{table}[t]
\centering
\caption{Preliminary monolingual EER (\%) on benchmark v0 (527
read-speech speakers; 12{,}000 trials per language). Fine-tuned rows
are $\mathrm{mean} \pm \mathrm{std}$ over three seeds. Preliminary,
benchmark v0: read speech, near single-session, closed-set speakers
for the fine-tuned rows; v1 (SLCeleb) provides the open-set,
in-the-wild condition.}
\label{tab:prelim}
\begin{tabular}{lcc}
\toprule
\textbf{System} & \textbf{Si} & \textbf{Ta} \\
\midrule
Zero-shot SpeechBrain ECAPA-TDNN            & 4.78 & 3.21 \\
Zero-shot ReDimNet-B6                       & 2.71 & 1.48 \\
ReDimNet-B6 + AS-Norm (training-free)       & 2.07 & 0.90 \\
WavLM-base+ LoRA fine-tune (3 seeds)        & 3.73 $\pm$ 0.49 & 3.11 $\pm$ 0.61 \\
WavLM-base+ full fine-tune (3 seeds)        & \textbf{1.69 $\pm$ 0.06} & \textbf{1.41 $\pm$ 0.10} \\
\bottomrule
\end{tabular}
\end{table}

Table~\ref{tab:prelim} gives preliminary monolingual context from
benchmark v0. Deployment-relevant preliminary findings: full
fine-tuning beat LoRA at every seed on this closed-set benchmark
(contrary to the published PEFT-under-shift prediction; the v1
open-set rerun arbitrates), and per-language score calibration is
mandatory --- applying a calibrator fitted on one language to the
other inflates Cllr by up to $5.7\times$. Cross-lingual (Cs) trials
remain pending on benchmark v1's bilingual speakers.

\subsection{Headline results}
\label{sec:exp:headline}

\begin{table}[t]
\centering
\caption{Speaker-verification performance on the SL-SPV evaluation set.
Numbers are pooled / Sinhala-only / Tamil-only / cross-lingual EER (\%),
$\mathrm{mean} \pm \mathrm{std}$ over three random seeds.}
\label{tab:headline}
\begin{tabular}{lcccc}
\toprule
\textbf{Config} & \textbf{Pooled} & \textbf{Si} & \textbf{Ta} & \textbf{Cs} \\
\midrule
P0 from scratch    & [X.X $\pm$ Y.Y] & [X.X $\pm$ Y.Y] & [X.X $\pm$ Y.Y] & [X.X $\pm$ Y.Y] \\
P1 + fine-tune     & [X.X $\pm$ Y.Y] & [X.X $\pm$ Y.Y] & [X.X $\pm$ Y.Y] & [X.X $\pm$ Y.Y] \\
P2 full stack      & \textbf{[X.X $\pm$ Y.Y]} & \textbf{[X.X $\pm$ Y.Y]}
                   & \textbf{[X.X $\pm$ Y.Y]} & \textbf{[X.X $\pm$ Y.Y]} \\
\bottomrule
\end{tabular}
\end{table}

\subsection{Ablation analysis}
\label{sec:exp:ablation}

\begin{table}[t]
\centering
\caption{Single-feature ablations from the full stack (P2), seed 42.
$\Delta$ = pooled-EER increase when the named feature is removed.}
\label{tab:ablation}
\begin{tabular}{lcc}
\toprule
\textbf{Ablation} & \textbf{Pooled EER \%} & $\boldsymbol{\Delta}$ \\
\midrule
P2 (all features)   & [X.X] & 0 \\
\midrule
$-$ fine-tune       & [X.X] & $+$[Y.Y] \\
$-$ LLRD            & [X.X] & $+$[Y.Y] \\
$-$ lang-aux        & [X.X] & $+$[Y.Y] \\
$-$ AS-Norm         & [X.X] & $+$[Y.Y] \\
$-$ PLDA            & [X.X] & $+$[Y.Y] \\
\bottomrule
\end{tabular}
\end{table}

\section{Discussion}
\label{sec:discussion}

\subsection{Implications for biometric deployment}
\label{sec:disc:deploy}

Voice biometric deployment in bilingual environments faces three
operating points that this benchmark explicitly characterises:
monolingual enrolment matching the test-time language (Si / Ta
columns of Table~\ref{tab:headline}), cross-lingual enrolment-test
mismatch (Cs column), and pooled (deployment-time language unknown).
A smaller Si--Ta gap indicates a more deployable model when the
language is not known in advance; a smaller Cs--Pooled gap indicates
robustness to language switching between enrolment and verification.

\subsection{Threats to validity}
\label{sec:disc:threats}

[NOTE: Discuss after results land. Likely topics: benchmark v0's
read-speech, near single-session character and closed-set fine-tuning
evaluation (optimistic absolute EERs; v1 SLCeleb provides the
open-set, multi-session condition), single-corpus-per-language
generalisation, channel homogeneity, and possible label noise in the
multi-task language-ID lookup.]

\subsection{Future work}
\label{sec:disc:future}

The toolkit ships infrastructure for two avenues we did not exercise
in this paper: (i) self-supervised continued pretraining on
unlabelled SL podcast/broadcast audio (FEATURE-009, design-only)
to further close the cross-lingual gap; and (ii) DANN-style adversarial
language and channel invariance (FEATURE-008), particularly relevant
for telephony deployment where the recording channel differs from
the training distribution.

\section{Conclusion}
\label{sec:conclusion}

We presented the first reproducible speaker-verification benchmark for
Sri Lankan Sinhala and Tamil, evaluated on a two-tier protocol
(benchmark v0: 527 read-speech speakers from OpenSLR SLR52/SLR65;
benchmark v1: SLCeleb, our 280-speaker in-the-wild corpus, pending). Single-feature ablations isolate the contribution
of each lever in the full-stack configuration: cross-lingual fine-tune,
layer-wise LR decay, multi-task language-ID, AS-Norm, and PLDA. The
open-source toolkit, corpus-preparation pipeline, and reproducibility
checklist are released for downstream work on Sri Lankan biometrics.

\section*{Acknowledgements}
[Funding, computational resources, advisor acknowledgements.]

\bibliographystyle{plain}
\bibliography{../references}

\section*{Biographical notes}
\noindent
\textbf{[First Author]} is [position] at [institution]. Their research
interests include [topics]. [One-paragraph bio per author per IJBM
submission guidelines.]

\end{document}
